\documentclass[11pt]{article}
\usepackage{mathblatt}
\begin{document}
\blattkopf*{Basisaufgaben}{BAS-Z5}{Basisaufgaben · Zettel · BAS-Z5}
\einheitenkopf[][Zettel 5]{Basisaufgaben · Zettel 5}
\zweigzeile{ohne Taschenrechner · je Aufgabe 1 Punkt · Lösungen auf der Rückseite}
% \small: zehn Aufgaben mit bis zu zwei Koordinatensystemen auf einer Seite (Render 28.09.)
\small

% 1: Bruchteil einer Fläche bestimmen – brueche-dezimalzahlen-basis-k1-v4 (Original 2026-FOR-B1b)
\begin{aufgabe}{Bei welcher Figur ist genau $\frac{1}{2}$ grau? \hfill (P26F) \\ \kreuz{P} \kreuz{Q} \kreuz{R} \kreuz{S}}

P \kreissektor[0.8]{150}{} \quad Q \bruchrechteck[2.5]{3}{6} \quad R \bruchkreis[0.8]{2}{5} \quad S \bruchrechteck[2.5]{2}{6}
\end{aufgabe}

% 2: Bruchteil einer Größe berechnen – brueche-dezimalzahlen-basis-k2-v2 (Original 2018-OS-B1a)
\begin{aufgabe}{$\frac{3}{4}$ von $2$ m – wie viel Meter? \hfill (P18) \\ \leerfeld[m]}
\end{aufgabe}

% 3: Zahlen in verschiedenen Darstellungen vergleichen – brueche-dezimalzahlen-basis-k4-v4 (Original 2023-OS-B1f)
\begin{aufgabe}{Welcher Wert ist der kleinste: $0{,}8$; $8\,\%$; $0{,}8^2$; $\frac{3}{4}$? \hfill (P23) \\ \leerfeld}
\end{aufgabe}

% 4: Termwert berechnen – rationale-zahlen-basis-k1-v3 (Original 2026-FOR-B1g)
\begin{aufgabe}{Berechne den Wert von $7 \cdot (x - 1)$ für $x = -2$. \hfill (P26F) \\ \leerfeld}
\end{aufgabe}

% 5: Zahl zu Bedingung angeben – rationale-zahlen-basis-k3-v1 (Original 2015-OS-B1b)
\begin{aufgabe}{Gib eine Zahl an, die größer ist als $-75$. \hfill (P15) \\ \leerfeld}
\end{aufgabe}

% 6: Zehnerpotenzschreibweise umwandeln – potenzen-wurzeln-basis-k3-v2 (Original 2025-OS-B1d)
\begin{aufgabe}{$3\,900\,000 = 3{,}9 \cdot 10^{\square}$ – welche Hochzahl? \hfill (P25) \\ \leerfeld}
\end{aufgabe}

% 7: Prozentsatz berechnen – prozentrechnung-basis-k3-v1 (Original 2014-OS-B1e)
\begin{aufgabe}{Ein Sparguthaben von $5\,000\,€$ bringt in einem Jahr $150\,€$ Zinsen. Wie hoch ist der Zinssatz? \hfill (P14) \\ \leerfeld[\%]}
\end{aufgabe}

% 8: Lineare Gleichung lösen – lineare-gleichungen-basis-k1-v3 (Original 2024-OS-B1d)
\begin{aufgabe}{Löse die Gleichung $4 \cdot (x + 3) = 8$. \hfill (P24) \\ x = \leerfeld}
\end{aufgabe}

% 9: Punktprobe durchführen – lineare-funktionen-basis-k3-v1 (Original 2023-OS-B1i)
\begin{aufgabe}{Gegeben ist $f(x) = 2x - 1$. Welcher Punkt liegt nicht auf der Geraden? Kreuze an. \hfill (P23) \\ \kreuz{$(1|1)$} \kreuz{$(4|7)$} \kreuz{$(-2|-3)$}}
\end{aufgabe}

% 10: Lage eines Punktes zu den Achsen erkennen – symmetrie-abbildungen-basis-k2-v1 (Original 2020-OS-B1b)
\begin{aufgabe}{Genau einer der vier Punkte liegt auf der x-Achse. Welcher? \hfill (P20) \\ \kreuz{$P(5|6)$} \kreuz{$Q(0|6)$} \kreuz{$R(5|0)$} \kreuz{$S(-5|-6)$}}
\end{aufgabe}

\begleitteil
\erg{1}{Q, denn $\frac{3}{6} = \frac{1}{2}$}
\erg{2}{$2 : 4 \cdot 3 = 1{,}5$ m}
\erg{3}{$8\,\% = 0{,}08$; die anderen sind $0{,}8$, $0{,}64$ und $0{,}75$}
\erg{4}{$7 \cdot (-2 - 1) = 7 \cdot (-3) = -21$}
\erg{5}{z. B. $-70$; richtig ist jede Zahl rechts von $-75$, auch $0$, nicht $-80$}
\erg{6}{$6$ – das Komma wandert sechs Stellen nach links}
\erg{7}{$150 : 5\,000 = 0{,}03 = 3\,\%$}
\erg{8}{$x + 3 = 2$, also $x = -1$}
\erg{9}{$(-2|-3)$, denn $f(-2) = -5$}
\erg{10}{$R(5|0)$}
\end{document}
